\section{Limitações e continuidade verificável}
A primeira limitação é o corpus parcial em termos materiais: foi obtida leitura textual do GDD, mas não sua cópia binária para hash, conferência de páginas e inspeção visual. A versão interna divergente do nome do arquivo também exige confirmação humana, sem presumir equivalência entre todas as versões mencionadas.

A segunda limitação é a inspeção estática: os resultados indicam estruturas e condições no código, mas não demonstram sua manifestação em runtime. Os arquivos de teste existentes no repositório são artefatos consultáveis, e sua presença não comprova execução recente, êxito em plataformas específicas ou cobertura das seis unidades.

A terceira limitação é analítica: seis unidades selecionadas não representam todo o GDD, e as categorias ainda não foram revisadas por pesquisadores humanos. O instrumento não possui avaliação de concordância e não sustenta estimativas de confiabilidade entre avaliadores; qualquer medida futura deverá informar procedimento, unidades, denominador e divergências observadas.

A quarta limitação é de proveniência: faltam registros originais que identifiquem quem decidiu cada parâmetro, o que impede separar contribuições individuais de alunos, orientador e agente. A assistência de IA à própria análise também deve ser considerada na revisão humana, especialmente na escolha de unidades, interpretação de ambiguidades e redação das conclusões.

A continuidade proposta consiste em conferir os bytes e a identidade da fonte, revisar humanamente as seis classificações, ampliar a extração às demais regras do GDD e executar verificações técnicas com versão do ambiente e resultados registrados. Essas etapas deverão conservar a distinção entre evidência documental, execução técnica e eventual avaliação educacional, para que mudanças de escopo não sejam apresentadas como resultados do piloto atual.
